% ============================================================================
% 01_intro.tex — merged paper (auction-v2)
% Nine-paragraph skeleton per plan/plan.md §2, re-sequenced per PI directive:
% auctions = main domain, DA/matching demoted to appendix robustness; menu per
% F1; clock-framing cross-model per F2; baseline provenance declared per F3;
% OSP-DA rescoped per F4; tier structure per F5 (Kendall's W + honest menu
% boundary); frontier nuance (F6) and the Section-7 cardinal question (F9)
% folded into P8a; direction shares use the repo-canonical convention (F9).
% Sources: writeup/contents/intro.tex; Engineering_simplicity/main.tex:69–89;
% results/merged_ranking/{auction_cells_summary.md, da_cells_summary.md,
% concordance.md}; results/cardinal/cardinal_summary.md;
% results/traces/traces_summary.md.
% ============================================================================
\section{Introduction}\label{sec:intro}

% --- P1: The design question -----------------------------------------------
A mechanism can be incentive compatible on paper and still fail in the minds of its participants. Truthful bidding is the dominant strategy in the second-price sealed-bid (SPSB) auction \citep{vickrey1961counterspeculation}, and the logic is short: a bid determines only \emph{whether} the bidder wins, not \emph{what} she pays; overbidding risks paying more than value, underbidding risks losing profitable wins. Yet human subjects have failed to bid truthfully in this auction for decades \citep{kagel1993independent, kagel1995individual}. The theory of simple mechanisms was built to explain such failures: obvious strategy-proofness \citep{li2017obviously}, strategic simplicity \citep{borgers2019strategically}, and one-step simplicity \citep{pycia2023theory} characterize which \emph{implementations} of the same allocation rule make good play cognitively feasible, and \citet{li2024designing} decomposes the underlying cognitive burden into contingent reasoning, forward planning, and belief formation. As large language model (LLM) agents begin to transact---bidding in ad exchanges, negotiating on procurement platforms, submitting preference rankings on behalf of applicants---this body of theory acquires a new and operational audience. The question is no longer only which designs are simple for humans, but which of the levers available to a mechanism designer preserve dominant-strategy play for a new kind of bounded reasoner. The stakes of the question are visible in a single cell of our data: GPT-4o's average bid in an SPSB auction falls roughly \$3.3 below its value, against dominant-strategy bids averaging \$24.5\footnote{Baseline provenance, here and throughout: comparisons across mechanisms (e.g., sealed bid versus clock) use the dedicated SPSB baseline cells (mean deviations $-\$0.49$, $-\$1.63$, $-\$3.28$, $-\$5.27$ for Claude, Gemini, GPT-4o, and Gemma respectively), while presentation and scaffold contrasts use the pooled axis-1/axis-3 baseline cells run in the same grid (a provenance audit reclassified the legacy axis-2 ``baseline'' as a treated clock-framing description; Section~\ref{sec:ranking-synthesis}); both conventions are reported in the master cell dataset (Section~\ref{sec:design}).}---yet a one-sentence change in how the very same rule is \emph{described} removes roughly half of that deviation.

% --- P2: What we do ----------------------------------------------------------
This paper delivers an empirical ranking of those levers. We organize the design space into three families: changes to the \emph{extensive form} (replacing a sealed bid with an ascending clock; replacing a one-shot report with an iterative query protocol), changes to the \emph{description} of a fixed mechanism (menu restatements, safety-exposing descriptions, clock-framing), and \emph{reasoning scaffolds} appended to the agent's instructions (contingent reasoning, forward planning, belief formation); a fourth family of preference framings is studied in Appendix~\ref{app:prospect}. Our main domain is the auction laboratory: an SPSB auction with three bidders and independent private values, surrounded by the first-price, third-price, common-value, and clock formats that make cross-format comparisons possible. As a generalization check, we replicate the full lever grid in a second canonical strategy-proof domain---student-proposing deferred acceptance (DA) in a $4\times4$ school-choice market with Ergin-acyclic priorities---reported in Appendix~\ref{app:da} and drawn on below where it sharpens the auction findings. We run the grid on four model families---GPT-4o (2024-08-06), Claude 3.5 Haiku (2024-10-22), Gemini 2.0 Flash, and Gemma 27B---and summarize each lever by a \emph{preservation index} $\rho$, the fraction of the baseline deviation from dominant-strategy play that the lever removes, computed within model and domain. We anchor the exposition on GPT-4o, which remains a widely used, well-studied general-purpose model for strategic and economic reasoning tasks in the recent literature \citep{horton2023large, manning2024automated}. This is a design choice for interpretability rather than a claim about model superiority: a fixed anchor isolates variation induced by mechanism, information environment, and presentation, while the ranking itself is computed within each model, so it is invariant to cross-model level differences.\footnote{Appendix~\ref{app:ablations} brackets the capability range. A frontier reasoning model (gpt-5-mini) bids the dominant strategy almost perfectly in SPSB (SMAD 0.2\%) and shades the first-price auction almost exactly to the risk-neutral equilibrium (mean deviation from the equilibrium benchmark $+\$0.002$; $-\$8.53$ from value), yet misses the third-price equilibrium badly (SMAD 60\%)---frontier capability narrows, but does not close, the bounded-rationality regime our ranking speaks to. A small open-weight model (Llama-3-8B, SMAD 57.58\%) marks the capability floor.}

% --- P3: Why the testbed is credible ----------------------------------------
A ranking computed on LLM agents is only as informative as the testbed that produces it, and a central challenge in using LLMs for mechanism design is deciding what they should be benchmarked against. Comparing bids only to theoretical equilibrium is informative but incomplete: mechanism performance in practice depends on systematic and well-documented deviations in human behavior. Auctions are a particularly useful proving ground precisely because they combine sharp theoretical predictions---equilibrium bidding and dominant-strategy truth-telling---with a deep experimental literature documenting robust regularities and mistakes, including overbidding, winner's-curse losses, and last-minute bidding under hard-close rules \citep{kagel2020handbook}. We therefore benchmark LLM behavior against human auction data as well as against theory. Using method-of-moments reconstructions of human bidding strategies from classic experiments \citep{kagel1993independent, kagel1986winner, li2017obviously, breitmoser2022obviousness, gonczarowski2024describing} and a unified metric---the scaled mean absolute deviation (SMAD)---we show in Section~\ref{sec:validation} that the testbed reproduces the human difficulty ranking of auction formats: GPT-4o's SMAD is 27.55\% in first-price versus 8.69\% in second-price auctions, mirroring the human 24.76\% versus 5.65\%. LLM bidders also reproduce the winner's-curse comparative statics---mean winner profits in first-price common-value auctions fall from $-2.54$ with four bidders to $-5.55$ with seven, and the curse is attenuated under second-price rules---and they snipe under hard-close rules in an eBay-style environment, echoing field evidence \citep{roth2002last}. Behavior is stable across rounds, with no evidence of within-session learning, so these patterns are properties of the agents rather than artifacts of experience.

% --- P4: Ranking result 1 — OSP tops (auctions-led; DA one sentence) ----------
The first ranking result is that obviously strategy-proof (OSP) extensive forms sit alone at the top. Replacing the sealed bid with an ascending clock compresses the cross-model mean deviation from $-\$2.67$ to roughly $-\$0.5$ (cross-model means, here and throughout, are unweighted averages over the four models), cuts SMAD to 1.4--5.2\% from baselines of 9.5--20.7\% ($p<0.001$ in every model, effect sizes up to $d=1.00$), and eliminates the heavy left tail of severe underbids.\footnote{A dedicated clean-IPV clock cell (IPV description and draws, $\$0.5$ tick, $K=50$) reproduces the result exactly for the two models re-runnable at the time of writing---SMAD $1.5\%$ (GPT-4o) and $1.8\%$ (Gemma)---so the rung does not rest on the legacy cells' mislabeled environment (Section~\ref{sec:ranking-extensive}).} The pattern generalizes beyond auctions. Standard DA is strategy-proof for the proposing side \citep{roth1982economics} but generally not OSP-implementable \citep{ashlagi2018stable}; under Ergin-acyclic priorities \citep{ergin2002efficient}, an iterative query protocol is OSP, and in our matching domain it drives observable misreports to zero for three of four models in the pick-based protocol (rule-of-three 95\% upper bounds below 1\%)---to our knowledge the first experimental test of an OSP implementation of DA with any kind of agent---while a yes/no implementation of the same protocol exposes model-specific failures that the pick-based accounting cannot see (Appendix~\ref{app:da}).

% --- P5: Ranking result 2 — descriptions split by invariance content ----------
The second result is that descriptions split into two sharply different sub-families, and the split is governed by what the words \emph{contain}, not by how many of them there are. A menu restatement in the style of \citet{gonczarowski2023strategyproofness}---rewording the SPSB rule as a menu of outcomes determined by others' bids---has no consistent effect: across the four models the shift is small and sign-mixed (Gemma moves toward truthful play, $+\$2.03$, $p<0.001$; GPT-4o away, $-\$1.16$, Welch $p=0.013$; Claude marginally away; Gemini is null), and null on average (placebo sign test $p=0.73$)---consistent, on average, with the modest behavioral effects of menu descriptions in human experiments \citep{gonczarowski2024describing}, though with model heterogeneity humans were never tested for. Safety-exposing descriptions, by contrast, are powerful: stating the payoff-safety invariant of the auction---your bid determines \emph{whether} you win, not \emph{what} you pay---cuts the cross-model mean deviation from $-\$2.67$ to $-\$1.53$. Clock-framing---merely \emph{describing} the sealed bid as an automatic exit threshold on a price clock, a change of words with no change to the game---helps three of the four models, cutting GPT-4o's mean deviation from $-\$2.94$ to $-\$1.54$ and halving Gemma's ($p<0.001$ each), while pushing the fourth, Gemini, away from truthful play---a first sign of a paraphrase sensitivity that Section~\ref{sec:ranking-descriptions} measures directly. The matching domain sharpens the diagnosis (Appendix~\ref{app:da}): there, a menu description that carries the invariance statement behaves like a safety description (pooled misreport error falls from 4.2\% to 1.8\%), a menu that restates only the computational mechanics makes play \emph{worse} (4.2\% to 6.9\%), and stating the rejection-safety invariant collapses error to 0.2\%, nearly matching the full OSP implementation. The active ingredient is not rewording the rules but exposing the incentive invariant that makes truth-telling safe.

% --- P6: Ranking result 3 — scaffolds are double-edged ------------------------
The third result is that reasoning scaffolds are double-edged, and the direction of the edge is predicted by which cognitive dimension they target. Scaffolding \emph{contingent reasoning}---a payoff-tree prompt that walks the agent through what happens in each outcome it might face---improves play across the board, cutting auction deviations from $-\$2.67$ to $-\$1.13$. Scaffolding \emph{belief formation} backfires: prompting first- and second-order beliefs about opponents pushes auction deviations to $-\$3.04$ and $-\$3.47$, with the damage concentrated in the weakest models while Gemini is largely immune. A one-shot sealed bid has no forward-planning axis, so \emph{forward planning} is tested in the matching domain (Appendix~\ref{app:da}), where it backfires as well: prompting models to reason about their own future decision points in DA's rejection dynamics worsens misreport error from 4.2\% to 7.8\% under one-step lookahead. The scaffolds give models material to strategize with, but not the ability to strategize well.

% --- P7: The assembled ranking, its concordance, and its theory parallel ------
Assembled, these results form the paper's central object: an empirical simplicity-preservation hierarchy of design levers (Section~\ref{sec:ranking}, Figure~\ref{fig:ranking}, Table~\ref{tab:ranking}). The hierarchy is a \emph{tier structure}, not a strict total order: OSP extensive forms at the top; safety-exposing descriptions, clock-framing, and contingent-reasoning scaffolds in a strong second tier whose internal order flips across domains (the payoff tree beats the safety description in three of four auction cells; the safety description beats the tree in all four matching cells); menu restatements statistically indistinguishable from baseline on average; forward-planning and belief scaffolds below baseline; and risk-averse preference framings at the bottom (Appendix~\ref{app:prospect}). The ordering is concordant across models and domains---Kendall's $W=0.70$ in auctions ($p=0.004$), $0.90$ in matching ($p=4\times10^{-4}$), and $0.81$ pooled ($p=8.4\times10^{-7}$)---and sign tests separate every adjacent tier boundary except one, which we report as an honest boundary failure: restatements do not reliably beat the strategizing scaffolds below them (4 of 8 cells, $p=1.0$), because the menu rung is sign-mixed rather than a clean placebo. That exception sharpens rather than muddies the lesson: pooling the evidence from both domains, the operative distinction within descriptions is \emph{invariance content}---safety-exposing text and the invariance-stating menu help, mechanical restatement is null to harmful. This localizes the binding constraint: what limits these agents is not \emph{computing} the dominant strategy---the payoff tree shows they can execute the case analysis when guided---but \emph{seeing that truth-telling is safe} without being shown. The models' own reasoning traces corroborate the mechanism with a complete dissociation---no lever moves stated reasoning and bids together: the payoff-safety description and the payoff tree move bids substantially while leaving stated reasoning unchanged (the safety invariant is echoed in none of the treated traces), whereas a worst-case contingency scaffold shifts stated reasoning sharply without moving bids (Section~\ref{sec:ranking}).

% --- P8: Human anchors, and the divergence that strengthens the claim ---------
Where human experimental evidence exists, it agrees with the ranking in sign on every rung: humans, too, improve under ascending clocks \citep{li2017obviously}---human preservation indices $+0.62$ (open clock) and $+0.37$ (blind clock)---improve under clock-framing \citep{breitmoser2022obviousness}---as do three of our four families---and are essentially unmoved by menu restatements \citep{gonczarowski2024describing}, whose human index of $-0.07$ matches the LLM null on average even as individual models diverge in sign. The agreement is more informative than it first appears, because the two populations make \emph{opposite} baseline mistakes: humans overbid in SPSB auctions---67.2\% of bids strictly exceed value, or 92.2\% among bids that deviate at all---while GPT-4o overbids in only 0.4\% of baseline bids and errs instead by shading down (Section~\ref{sec:validation}).\footnote{Direction shares use the repository-canonical convention: the share of all bids strictly above value, with the share among deviating bids reported alongside. Earlier drafts reported a 96.0\% human overbid share under a $\pm2\%$ truth-telling-band convention.} That the same levers repair opposite errors in both populations indicates that the levers operate on the strategic problem itself---the recognizability of the dominant strategy---rather than on any psychology particular to humans. It also disciplines interpretation: our LLM testbed is a model of a boundedly rational player, not a simulation of a human subject. The rungs with only partial human anchors are then delivered as predictions (Section~\ref{sec:humans}): most cheaply, that a safety-exposing description---an honest, one-sentence statement of the mechanism's incentive invariant---will shift human bidding toward value in auctions, where, unlike matching \citep{gonczarowski2024describing, katuscak2024simpler}, no invariant-statement experiment yet exists.

% --- P8a: Scope — frontier capability and the ordinal/cardinal boundary -------
Two boundaries discipline the scope of these findings, and the paper closes by measuring both. The first is capability. Frontier reasoning does not uniformly dissolve the constraints: gpt-5-mini plays the dominant strategy almost perfectly in SPSB and applies the textbook equilibrium-shading formula in first price, yet misses the third-price equilibrium badly---it dissolves the constraints exactly where optimal play is textbook, and nowhere else (Appendix~\ref{app:ablations}). Nor does frontier capability immunize against presentation: in a four-model frontier intervention grid, a content-free menu restatement collapses one otherwise exactly-truthful model to severe underbidding (Section~\ref{sec:ranking-synthesis}). The second is the depth of the human agreement. Everything above establishes \emph{ordinal} agreement: LLM populations rank formats and levers the way humans do. Section~\ref{sec:humans} then asks what it would take to agree \emph{cardinally}---whether any tuning of the agents (temperature, risk personas, preference framings, rule explanations) makes their bidding match human behavior in levels and direction. The answer is that no single tuning does: the best level-matching knob is mechanism-specific, and the only knob that reproduces human overbidding in second- and third-price auctions---a risk-seeking persona---simultaneously introduces dominated overbidding in the first-price auction, where humans almost never bid above value. Ordinal fidelity is a property of the agent population; cardinal fidelity is, at best, a mechanism-by-mechanism calibration.

% --- P9: Implications and contributions ---------------------------------------
The paper makes four contributions. First, the ranking itself: a tiered, cross-model ordering of the mechanism designer's levers by how much dominant-strategy play they preserve, with a common preservation index making rungs comparable, established in auctions and replicated in matching. Second, the first experimental test of an OSP implementation of deferred acceptance with any kind of agent, with decision-level error accounting that makes precise what ``zero observable misreports'' does and does not cover (Appendix~\ref{app:da}). Third, a methodology: moment-matched reconstruction of human experimental benchmarks and the SMAD metric family, which turn decades of laboratory results into a validation layer for agentic testbeds---together with a characterization of where that validation stops, at the ordinal--cardinal boundary\ifanon\else; we release the framework as an open-source benchmark\fi. Fourth, design guidance. For a designer who must keep a direct mechanism, an honest safety-exposing description recovers most of the benefit of full OSP redesign at the cost of a sentence; conversely, prompting agents to plan ahead or model others' beliefs---precisely the ``reasoning enhancements'' an agent vendor might advertise---systematically backfires. If the cognitive constraints that the theory of simplicity was built to address also bind on artificial agents, then the case for engineering simplicity no longer rests on human psychology alone: it becomes a design principle for markets that must work for any kind of reasoning agent.

% --- P10: Roadmap --------------------------------------------------------------
The paper proceeds as follows. Section~\ref{sec:related} reviews the related literature. Section~\ref{sec:framework} derives the lever taxonomy from simplicity theory and states our hypotheses. Section~\ref{sec:design} describes the mechanisms, environments, treatments, models, metrics, and the reconstruction of human benchmarks. Section~\ref{sec:validation} establishes the testbed's behavioral fidelity against human experimental and field evidence. Section~\ref{sec:ranking} presents the ranking in the auction domain: extensive forms, then descriptions, then scaffolds, then the assembled hierarchy, together with the reasoning-trace evidence on how the levers act. Section~\ref{sec:humans} moves from ordinal to cardinal agreement with human behavior: it anchors the ranking rung by rung to human evidence, issues out-of-sample predictions for unanchored rungs, and asks whether any tuning of the agents achieves cardinal agreement across mechanisms. Section~\ref{sec:discussion} draws out design implications, scope conditions, and limitations, and Section~\ref{sec:conclusion} concludes. Appendices collect the human-data reconstructions, simulation procedures, prompt library, additional ablations, and the full replication of the lever grid in the school-choice matching domain (Appendix~\ref{app:da}).
